\input{../tex-inputs/latex-leseansicht-vorspann}

\section[Arthur und Olga Schnitzler an Gustav Schwarzkopf, 18. 9. 1906]{L04387 Arthur und Olga Schnitzler an Gustav Schwarzkopf, 18. 9. 1906}
\nopagebreak\mylabel{L04387v}
\rehead{ }\normalsize\beginnumbering\briefempfaengerindex{Schwarzkopf, Gustav@\textsc{Schwarzkopf, Gustav}!zzzSchnitzler, Olga@\emph{von Olga Schnitzler}!1906-09-181@{18. 9. {[}1906{]}}|(be}\briefempfaengerindex{Schwarzkopf, Gustav@\textsc{Schwarzkopf, Gustav}!zzzSchnitzler, Arthur@\emph{von Arthur Schnitzler}!1906-09-181@{18. 9. {[}1906{]}}|(be}
\toendnotes[C]{\smallbreak\pagebreak[2]}
\correspDesc{Versand  durch Arthur Schnitzler, Olga Schnitzler am 18. 9. [1906] in Semmering
\newline{}Übermittlung  am 18. 9. 1906 in Semmering
\newline{}Erhalt  durch Gustav Schwarzkopf im Zeitraum [19. 9. 1906
                  – 23. 9. 1906?] in Tiefer Graben 23}\toendnotes[C]{\smallbreak}
\Standort{DLA, A:Schnitzler, HS.1985.1.1897.}
\physDesc{Bildpostkarte, 306 Zeichen
\newline{}Handschrift Arthur Schnitzler: Bleistift, deutsche Kurrent
\newline{}Handschrift Olga Schnitzler: Bleistift
\newline{}Versand: Stempel: »\nobreak{}\oindex{Semmering@\textbf{Semmering}|pwk}Semmering, 1\textcolor{gray}{8} IX 06, 3\nobreak{}«.  }\pstart{}{\pb}Hrn Gustav Schwarzkopf\pend{}\pstart{}Wien I\oindex{I., Innere Stadt@\textbf{I., Innere Stadt}|pw}\pend{}\pstart{}Tiefer Graben 23\oindex{Wien@\textbf{Wien}!I., Innere Stadt@\textbf{I., Innere Stadt}!Tiefer Graben 23@\textbf{Tiefer Graben 23}|pw}.\pend{}{\bigskip}
\pstart
           \centering{}{\pb}\textcolor{gray}{\textbf{Südbahnhotel\oindex{Südbahnhotel [Semmering]@\textbf{Südbahnhotel [Semmering]}|pw} u. Waldhof Semmering\oindex{Waldhof, Dependance des Südbahnhotels Semmering@\textbf{Waldhof, Dependance des Südbahnhotels Semmering}|pw}.}}\pend
           \vspace{1em}
\pstart
           \noindent{}{\pb}lieber Guſtav, wir ſind hier{ }ſeit 8 Tagen u finden, daſs auch
               das ſchlechte Wetter hier ſchöner iſt als in Wien\oindex{Wien@\textbf{Wien}|pw}.
               Ende der Woche hoffen wir wieder daheim zu{ }ſein u{ }ſehen Sie bald. Auch \textsc{Hugo}’s\pwindex{Hofmannsthal, Hugo von 1.\,2.\,1874 Wien – 15.\,7.\,1929 Rodaun@\textsc{Hofmannsthal, Hugo von} (1.\,2.\,1874 Wien – 15.\,7.\,1929 Rodaun)|pw}\pwindex{Hofmannsthal, Gertrude von 16.\,3.\,1880 Wien – 9.\,11.\,1959 Paddington@\textsc{Hofmannsthal, Gertrude von} (16.\,3.\,1880 Wien – 9.\,11.\,1959 Paddington)|pw} u \textsc{Fred\pwindex{W. Fred 29.\,6.\,1879 Wien – 23.\,10.\,1922 Berlin@\textsc{W. Fred} (29.\,6.\,1879 Wien – 23.\,10.\,1922 Berlin)|pw}}{ }ſind hier.\pend
           
\pstart
           Herzlichſt, mit Grüßen {\\[\baselineskip]}an Sie u Dr. Max\pwindex{Schwarzkopf, Max 12.\,6.\,1857 Wien – 14.\,4.\,1928 ebd.@\textsc{Schwarzkopf, Max} (12.\,6.\,1857 Wien – 14.\,4.\,1928 ebd.)|pw}{\\[\baselineskip]}Ihr \spacefill\mbox{Arthur}  \pend
           \leftskip=0em{}
\pstart
           18. 9.\pend
           \selectlanguage{ngerman}\vspace{1em}\pstart \spacefill\mbox{{[}hs. O. Schnitzler:{]} O. S.}\pend{}\selectlanguage{ngerman}\endnumbering\briefempfaengerindex{Schwarzkopf, Gustav@\textsc{Schwarzkopf, Gustav}!zzzSchnitzler, Olga@\emph{von Olga Schnitzler}!1906-09-181@{18. 9. {[}1906{]}}|)be}\briefempfaengerindex{Schwarzkopf, Gustav@\textsc{Schwarzkopf, Gustav}!zzzSchnitzler, Arthur@\emph{von Arthur Schnitzler}!1906-09-181@{18. 9. {[}1906{]}}|)be}\mylabel{L04387h}
\begin{anhang}
\end{anhang}\newcommand{\dateiname}{L04387}\newcommand{\titel}{Arthur und Olga Schnitzler an Gustav Schwarzkopf, 18. 9. 1906}\newcommand{\editorInnen}{Herausgegeben von Jahnke, SelmaMüller, Martin Anton}\input{../tex-inputs/latex-leseansicht-abspann}
